\begin{frame}{Four assignments, four rows}
\princ{EXHAUST SMALL CASES}
\renewcommand{\arraystretch}{1.35}
\begin{center}
\begin{tabular}{c c | c | c | c}
$P$ & $Q$ & $P\ \text{OR}\ Q$ & $P\ \text{OR}\ \text{NOT }Q$ & $F$ \\\midrule
False & False & False & True & False\\
False & True & True & False & False\\
True & False & True & True & True\\
True & True & True & True & True
\end{tabular}
\end{center}
\predict{For which row could $B=P\ \text{OR}\ Q$ disagree with $F$?}
\see{the row $P=False, Q=True$ makes $B=True$ while $F=False$.}
\note{Fill the row values with the class if possible. The point is not speed; it is covering the entire finite space.}
\end{frame}

\begin{frame}{One counterexample is enough}
\princ{COUNTEREXAMPLE}
\begin{center}
\begin{tabular}{c c | c | c | c}
$P$ & $Q$ & $F$ & $B$ & match?\\\midrule
False & False & False & False & yes\\
\row{False}{True}{False}{True}
True & False & True & True & yes\\
True & True & True & True & yes
\end{tabular}
\end{center}
\begin{alertblock}{Conclusion}
The proposed simplification is not equivalent to $F$ because one assignment differs.
\end{alertblock}
\see{the exact mismatch is $P=False, Q=True$; no larger example is needed.}
\end{frame}

\begin{frame}[fragile]{Prompt 01: make the claim checkable}
\princ{DEFINE BEFORE COMPUTING}
\promptbox{01}{a}
\predict{What must be defined before an evaluator can be trusted?}
\see{terms, operators, assignment order, and the meaning of equivalence are visible.}
\note{Have students identify a missing definition before running anything.}
\end{frame}

\begin{frame}[fragile]{Prompt 02: make every row visible}
\princ{EXHAUST SMALL CASES}
\promptbox{02}{a}
\predict{Which row should reveal the false simplification?}
\see{four rows and exactly one mismatch: $P=False, Q=True$.}
\end{frame}

\begin{frame}{A correct simplification from the activity}
\princ{COMPARE CLAIMS}
The Tuesday source activity supplies:
\[
E=(P\ \text{AND}\ Q)\ \text{OR}\ (P\ \text{AND}\ \text{NOT }Q)
\]
and proposes $S=P$.
\begin{itemize}
\item Complete the same four rows.
\item $E$ and $S$ agree on every row.
\item A correct simplification is earned by the table, not by familiarity.
\end{itemize}
\see{the source example and the deliberately false example use the same checking method.}
\end{frame}

\begin{frame}{From expression to circuit}
\princ{REPRESENTATION}
\begin{columns}[T]
\begin{column}{0.58\textwidth}
\begin{tikzpicture}[node distance=7mm, every node/.style={font=\small}, gate/.style={draw,rounded corners,fill=blue!10,minimum width=1.35cm,minimum height=0.65cm}, wire/.style={-{Stealth[length=4pt]},thick}]
\node (p) at (0.2,2.2) {$P$};
\node (q) at (0.7,0.55) {$Q$};
\node (or1) [gate] at (2,2.2) {OR};
\node (notq) [gate] at (2,0) {NOT $Q$};
\node (or2) [gate] at (4,0) {OR};
\node (and) [gate] at (6,1.1) {AND};
\node (out) at (7.4,1.1) {$F$};
\coordinate (psource) at (0.7,2.2);
\coordinate (qsource) at (1.2,0.55);
\coordinate (or1inupper) at ([yshift=3pt]or1.west);
\coordinate (or1inlower) at ([yshift=-3pt]or1.west);
\coordinate (or2inupper) at ([yshift=3pt]or2.west);
\coordinate (or2inlower) at ([yshift=-3pt]or2.west);
\coordinate (andinupper) at ([yshift=3pt]and.west);
\coordinate (andinlower) at ([yshift=-3pt]and.west);
\draw[wire] (p) -- (psource);
\draw[wire] (psource) -- (or1inupper);
\draw[wire] (psource) -- ++(-0.2,0) |- (0.5,-0.5) -- (3.15,-0.5) |- (or2inlower);
\draw[wire] (q) -- (qsource);
\draw[wire] (qsource) -- ++(-0.1,0) |- (or1inlower);
\draw[wire] (qsource) -- (notq.west);
\draw[wire] (notq.east) -- (or2inupper);
\draw[wire] (or1.east) -| (andinupper);
\draw[wire] (or2.east) |- (andinlower);
\draw[wire] (and.east) -- (out);
\end{tikzpicture}
\end{column}
\begin{column}{0.38\textwidth}
A circuit is a representation: each gate computes one operation and each wire carries a value.\par\medskip
\textbf{Check:} a circuit that omits a wire may implement a different expression.
\end{column}
\end{columns}
\predict{If one wire is missing, should we trust the same truth table?}
\see{the picture must be translated back into an expression before checking it.}
\end{frame}

\begin{frame}{Equivalence is a universal statement}
\princ{MATCH THE SCOPE}
\begin{block}{Equivalent}
Two expressions are equivalent when they return the same value for every assignment in the declared scope.
\end{block}
\begin{itemize}
\item “Every” is why the four-row table matters.
\item One mismatch changes an equivalence claim to “not equivalent.”
\item A matching sample is evidence for those rows, not a universal proof.
\end{itemize}
\see{scope = variables $P,Q$ and operators AND, OR, NOT in these examples.}
\end{frame}
